\subsection{下中心列，幂零群}

设 \(\mathbf{G}\) 为一个群， \(\mathbf{H}\) 为 \(\mathbf{G}\) 的一个子群， \(\mathbf{K}\) 为 \(\mathbf{G}\) 的一个正规子群。\(\mathbf{H}\) 在 \(\mathbf{G} / \mathbf{K}\) 中的像包含于 \(\mathbf{G} / \mathbf{K}\) 的中心，当且仅当 \((\mathbf{G}, \mathbf{H}) \subset \mathbf{K}\) 。

定义 5. 设 G 是一个群。G 的下中心列是 G 的子群序列 \((\mathbf{C}^n (\mathbf{G}))_{n\geqslant 1}\) ，归纳地定义如下：

\[
\mathbf {C} ^ {1} (\mathbf {G}) = \mathbf {G}, \quad \mathbf {C} ^ {n + 1} (\mathbf {G}) = (\mathbf {G}, \mathbf {C} ^ {n} (\mathbf {G})).
\]

设 \(f: \mathbf{G} \to \mathbf{G}'\) 是群同态。对 \(n\) 作归纳可以看出， \(f(\mathbf{C}^n(\mathbf{G})) \subset \mathbf{C}^n(\mathbf{G}')\) ，且当 \(f\) 是满射时， \(f(\mathbf{C}^n(\mathbf{G})) = \mathbf{C}^n(\mathbf{G}')\) 。特别地，对所有 \(n \geqslant 1\) , \(\mathbf{C}^2(\mathbf{G})\) 是 \(\mathbf{G}\) 的特征子群（从而是正规子群）。对所有 \(n \geqslant 1\) , \(\mathbf{C}^n(\mathbf{G}) / \mathbf{C}^{n+1}(\mathbf{G})\) 包含于 \(\mathbf{G} / \mathbf{C}^{n+1}(\mathbf{G})\) 的中心。

设 \((\mathbf{G}_1, \mathbf{G}_2, \ldots)\) 是 \(\mathbf{G}\) 的正规子群的一个递减序列，满足 (1) \(\mathbf{G}_1 = \mathbf{G}\) ；(2) 对所有 \(i\) , \(\mathbf{G}_i / \mathbf{G}_{i+1}\) 包含于 \(\mathbf{G} / \mathbf{G}_{i+1}\) 的中心。则 \(\mathbf{C}^i(\mathbf{G}) \subset \mathbf{G}_i\) ，这可对 \(i\) 作归纳看出。

现在

\[
\left(\mathbf {C} ^ {m} (\mathbf {G}), \mathbf {C} ^ {n} (\mathbf {G})\right) \subset \mathbf {C} ^ {m + n} (\mathbf {G}).\tag{7}
\]

因为，若此关系记为 \((\mathbf{F}_{m,n})\) ，则由 \((\mathbf{F}_{m,n})\) ，根据第2号命题5，可得

\[
\begin{array}{l} \left(\mathrm{C} ^ {m} (\mathrm{G}), \mathrm{C} ^ {n + 1} (\mathrm{G})\right) \subset \left(\mathrm{G}, \left(\mathrm{C} ^ {m} (\mathrm{G}), \mathrm{C} ^ {n} (\mathrm{G})\right)\right). \left(\mathrm{C} ^ {n} (\mathrm{G}), \left(\mathrm{G}, \mathrm{C} ^ {m} (\mathrm{G})\right)\right) \\ \subset \mathrm{C} ^ {m + n + 1} (\mathrm{G}). \left(\mathrm{C} ^ {m + 1} (\mathrm{G}), \mathrm{C} ^ {n} (\mathrm{G})\right). \end{array}
\]

因此 \(((\mathbf{F}_{m,n})\) 和 \((\mathbf{F}_{m + 1,n}))\Rightarrow (\mathbf{F}_{m,n + 1})\) 。由于 \((\mathbf{F}_{m,1})\) 与 \((\mathbf{F}_{1,n})\) 是显然的， \((\mathbf{F}_{m,n})\) 由归纳法可得。

定义 6. 群 G 称为幂零的，如果存在整数 n 使得 \(\mathbf{C}^{n+1}(\mathbf{G}) = \{e\}\) 。使 \(\mathbf{C}^{n+1}(\mathbf{G}) = \{e\}\) 的最小整数 n 称为幂零群 G 的幂零类。

若 \(n \in N\) ，幂零类为 n 的群称为 n 类幂零群。有时也说：若 G 幂零，则群 G 的幂零类是有限的。

例子。(1) 一个群是 0 类（相应地，类 \(\leqslant 1\) ）幂零的，当且仅当它由单位元组成（相应地，是交换的）。

\begin{enumerate}
\def\labelenumi{(\arabic{enumi})}
\setcounter{enumi}{1}
\item
  *对每个交换环 A 与每个整数 \(n \geqslant 1\) ，上严格三角群 \(\mathrm{T}_1(n, \mathbf{A})\) 的幂零类 \(\leqslant n - 1\) （且恰为类 \(n - 1\) ，若 \(\mathbf{A} \neq \{0\}\) ）。*
\item
  设 \(\mathbf{G}\) 为类 \(n\) 的幂零群。\(\mathbf{G}\) 的每个子群（相应地，每个商群）的幂零类都 \(\leqslant n\) 。事实上，若 \(\mathbf{H}\) 是 \(\mathbf{G}\) 的子群，则 \(\mathbf{C}^n (\mathbf{H})\subset \mathbf{C}^n (\mathbf{G})\) 。若 \(\mathbf{G}'\) 是 \(\mathbf{G}\) 的商群，且 \(\pi : \mathbf{G}\to \mathbf{G}'\) 是典范同态，则 \(\mathbf{C}^n (\mathbf{G}') = \pi (\mathbf{C}^n (\mathbf{G}))\) 。
\item
  幂零群的有限乘积是幂零群。
\end{enumerate}

命题 7. 设 G 是一个群，n 是一个整数。下列条件等价：

\begin{enumerate}
\def\labelenumi{(\alph{enumi})}
\item
  G 是类 \(\leqslant n\) 的幂零群。
\item
  存在 \(\mathbf{G}\) 的子群列：
\end{enumerate}

\[
\mathrm{G} = \mathrm{G} ^ {1} \supset \mathrm{G} ^ {2} \supset \dots \supset \mathrm{G} ^ {n + 1} = \{e \}
\]

使得 \((\mathbf{G},\mathbf{G}^k)\subset \mathbf{G}^{k + 1}\) 对所有 \(k\in [1,n]\) 成立。

\begin{enumerate}
\def\labelenumi{(\alph{enumi})}
\setcounter{enumi}{2}
\item
  存在子群 \(\mathbf{A}\) （ \(\mathbf{G}\) 的子群），包含在 \(\mathbf{G}\) 的中心内，使得 \(\mathbf{G} / \mathbf{A}\) 是类 \(\leqslant n - 1\) 的幂零群。
\item
  (a) \(\Rightarrow\) (b)：只需取 \(\mathbf{G}^k = \mathbf{C}^k (\mathbf{G})\) .
\item
  (b) \(\Rightarrow\) (a)：对 \(k\) 作归纳，得 \(\mathbf{C}^k (\mathbf{G})\subset \mathbf{G}^k\)
\item
  (a) \(\Rightarrow\) (c)：只需取 \(\mathbf{A} = \mathbf{C}^n (\mathbf{G})\) .
\item
  (c) \(\Rightarrow\) (a)：设 \(\pi: \mathbf{G} \to \mathbf{G} / \mathbf{A}\) 为典范同态；则 \(\pi(\mathbf{C}^n(\mathbf{G})) = \mathbf{C}^n(\mathbf{G} / \mathbf{A}) = \{e\}\) ，故 \(\mathbf{C}^n(\mathbf{G}) \subset \mathbf{A}\) ，从而 \(\mathbf{C}^{n+1}(\mathbf{G}) = \{e\}\) 。
\end{enumerate}

更简洁地说：一个群是类 \(\leqslant n\) 的幂零群，如果它可由群 \(\{e\}\) 经 n 次逐次中心扩张得到。

推论. 幂零群的中心扩张（被一个必然交换的群扩张）是幂零的。

命题 8. 设 G 是类 \(\leqslant n\) 的幂零群，H 是 G 的一个子群。存在子群序列

\[
\mathbf {G} = \mathbf {H} ^ {1} \supset \mathbf {H} ^ {2} \supset \dots \supset \mathbf {H} ^ {n + 1} = \mathbf {H},
\]

使得 \(\mathbf{H}^{k + 1}\) 在 \(\mathbf{H}^k\) 中正规，且 \(\mathbf{H}^k /\mathbf{H}^{k + 1}\) 是交换的（对所有 \(k\leqslant n\) ）

选取一个序列 \((\mathbf{G}^{k})\) 使得对每个 k 都满足命题 7(b) 的条件； \(G^{k}\) 在 G 中是正规的。记：

\[
\mathbf {H} ^ {k} = \mathbf {H}. \mathbf {G} ^ {k}.
\]

需要验证 \(\mathbf{H}^{k + 1}\) 被 \(\mathbf{H}^k = \mathbf{H}.\mathbf{G}^k\) 正规化；由于它被 \(\mathbf{H}\) 正规化，只需验证它被 \(\mathbf{G}^k\) 正规化即可。现在，若 \(s\in \mathbf{G}^k\) 与 \(h\in \mathbf{H}\) ，

\[
s h s ^ {- 1} = s h s ^ {- 1} h ^ {- 1}. h \in (\mathrm{G}, \mathrm{G} ^ {k}). \mathrm{H}
\]

而 \((\mathbf{G},\mathbf{G}^k)\) . H 包含于 \(\mathbf{G}^{k + 1}\) ，且 \(\mathbf{H} = \mathbf{H}^{k + 1}\) ；因此 \(s.\mathbf{H}^{k + 1}.s^{-1} = \mathbf{H}^{k + 1}\) ，这表明 \(\mathbf{H}^{k + 1}\) 在 \(\mathbf{H}^k\) 中正规。

最后，典范同态 \(G^{k}/G^{k+1} \to H^{k}/H^{k+1}\) 显然是满射；由于第一个群是交换的，第二个群也是交换的。

推论 1. 设 G 为幂零群，H 为 G 的子群。若 H 不同于 G，则 H 在 G 中的正规化子 \(N_{G}(H)\) 也不同于 H。

设 \(k\) 为使 \(H^{k} \neq H\) 的最大指标。群 \(H^{k}\) 正规化 \(H\) 且不同于 \(H\) 。

推论 2. 设 G 为幂零群，H 为 G 的子群。若 H 不同于 G，则存在 G 的正规子群 N，包含 H 且不同于 G，使得 G/N 为交换群。

设 \(k\) 是使 \(H^k \neq G\) 的最小指标。群 \(H^k\) 满足所要求的条件。

推论3. 设G为幂零群，H为G的子群。若G = H.(G, G)，则G = H。

G 的每个包含 H 且使 G/N 交换的子群 N 都包含 H.(G, G)。于是推论 3 由推论 2 得出。

推论 3 也可以这样表述：设 X 是 G 的一个子集。要使 X 生成 G，当且仅当 X 在 G/D(G) 中的像生成 G/D(G)。

推论 4. 设 \(f: \mathbf{G}' \to \mathbf{G}\) 是一个群同态。假设

\begin{enumerate}
\def\labelenumi{(\alph{enumi})}
\item
  G 是幂零的。
\item
  同态 \(f_1 \colon \mathbf{G}' / (\mathbf{G}', \mathbf{G}') \to \mathbf{G} / (\mathbf{G}, \mathbf{G})\) （由 \(f\) 过渡到商导出）是满射的。
\end{enumerate}

那么 \(f\) 是满射。

这由推论 3 应用于子群 \(\mathbf{H} = f(\mathbf{G}')\) 得出。

命题 9. 设 G 是类 \(\leqslant n\) 的幂零群，且 N 是 G 的一个正规子群。存在子群列

\[
\mathrm{N} = \mathrm{N} ^ {1} \supset \mathrm{N} ^ {2} \supset \dots \supset \mathrm{N} ^ {n + 1} = \{e \}
\]

使得 \((\mathbf{G},\mathbf{N}^k)\subset \mathbf{N}^{k + 1}\) 对 \(k = 1,\dots ,n\) 成立

若 \((\mathbf{G}^k)\) 满足命题 7 的条件 (b)，则取

\[
\mathbf {N} ^ {k} = \mathbf {G} ^ {k} \cap \mathbf {N}.
\]

推论 1. 设 G 为幂零群，Z 为 G 的中心，N 为 G 的正规子群。若 N ≠ \{e\}，则 N ∩ Z ≠ \{e\}。

设 \(k\) 为使 \(\mathbf{N}^k \neq \{e\}\) 的最大指标。群 \(\mathbf{N}^k\) 包含于 \(\mathbf{N}\) 。另一方面， \((\mathbf{G}, \mathbf{N}^k) \subset \mathbf{N}^{k+1} = \{e\}\) ；因此 \(\mathbf{N}^k\) 包含于 \(\mathbf{Z}\) （ \(\mathbf{G}\) 的中心）内。

推论 2. 设 \(f\) 是从幂零群 \(G\) 到群 \(G'\) 的一个同态。若 \(f\) 在 \(G\) 的中心上的限制是单射，则 \(f\) 是单射。

这就是推论 1 应用于 \(\operatorname{Ker}(f)\) 的情形。

